% =====================================================================
%  preamble.tex -- AMS-style variant.
%  amsart already supplies theorem environments (thm, lem, prop, cor,
%  defn, rem, conj, ...) and handles \title / \author / \date / abstract
%  itself, so this preamble is much shorter than the article one.
% =====================================================================

\usepackage[T1]{fontenc}
\usepackage[utf8]{inputenc}
\usepackage{lmodern}
\usepackage{amsmath,amssymb,amsthm,mathrsfs}
\usepackage{geometry}
\geometry{a4paper,margin=30mm}
\usepackage{microtype}
\usepackage{graphicx,booktabs,array}
\usepackage{enumitem}
\setlist{itemsep=3pt,topsep=5pt}
\usepackage{xcolor}
\definecolor{linkblue}{RGB}{31,70,105}
\usepackage{tikz}
\usetikzlibrary{arrows.meta,positioning,calc,decorations.pathmorphing,fit}

% hyperref must come after everything else that defines commands.
\usepackage[colorlinks=true,linkcolor=linkblue,citecolor=linkblue,
            urlcolor=linkblue]{hyperref}
\usepackage[capitalise,noabbrev,nameinlink]{cleveref}

\numberwithin{equation}{section}

% amsart numbers theorem-like environments continuously (1, 2, 3, ...).
% Uncomment the next line to number them within the section instead.
% \numberwithin{thm}{section}

% Convenience aliases so the AMS environment names read like the
% article-style names used elsewhere in this collection.
\newtheorem{theorem}{Theorem}
\newtheorem{lemma}[theorem]{Lemma}
\newtheorem{proposition}[theorem]{Proposition}
\newtheorem{corollary}[theorem]{Corollary}
\theoremstyle{definition}
\newtheorem{definition}[theorem]{Definition}
\newtheorem{example}[theorem]{Example}
\newtheorem{convention}[theorem]{Convention}
\theoremstyle{remark}
\newtheorem{remark}[theorem]{Remark}
\newtheorem{notation}[theorem]{Notation}

\allowdisplaybreaks[2]

% Project macros.
\newcommand{\N}{\mathbb{N}}
\newcommand{\Z}{\mathbb{Z}}
\newcommand{\Q}{\mathbb{Q}}
\newcommand{\R}{\mathbb{R}}
\newcommand{\C}{\mathbb{C}}
\newcommand{\F}{\mathbb{F}}
\DeclareMathOperator{\rank}{rank}
\DeclareMathOperator{\supp}{supp}
\DeclareMathOperator{\im}{im}
\DeclareMathOperator{\id}{id}
\DeclareMathOperator{\tr}{tr}
\DeclareMathOperator{\Spec}{Spec}
\DeclareMathOperator{\Hom}{Hom}
\DeclareMathOperator{\End}{End}
\newcommand{\dd}{\mathrm{d}}
\newcommand{\abs}[1]{\left\lvert #1 \right\rvert}
\newcommand{\norm}[1]{\left\lVert #1 \right\rVert}
\newcommand{\ip}[2]{\left\langle #1, #2 \right\rangle}
